% Hoja de fórmulas — Química General A1095 (UTN FRBA, Ing. Mecánica).
% Una sola carilla A4. Solo lleva las fórmulas que pidió el usuario; las
% sugerencias extra van únicamente en hoja_formulas.md.
% Compilar con: make  ->  build/hoja_formulas.pdf

\documentclass[9pt]{extarticle}
\usepackage[a4paper, margin=7mm]{geometry}
\usepackage{amsmath}
\usepackage{fontspec}
\usepackage{unicode-math}
\setmainfont{Latin Modern Roman}
\setmathfont{Latin Modern Math}
\usepackage{graphicx}
\usepackage{multicol}
\usepackage[dvipsnames]{xcolor}

\pagestyle{empty}
\frenchspacing
\setlength{\parindent}{0pt}
\setlength{\parskip}{0pt}
\setlength{\columnsep}{4mm}
\setlength{\columnseprule}{0.3pt}
\def\columnseprulecolor{\color{gray!50}}
\setlength{\abovedisplayskip}{1pt}
\setlength{\belowdisplayskip}{1pt}

% Capítulo: barra de color a todo el ancho de la columna.
\newcommand{\capitulo}[1]{%
  \par\vspace{2pt}\colorbox{NavyBlue}{\parbox{\dimexpr\linewidth-2\fboxsep}{%
  \color{white}\bfseries\small #1}}\par\vspace{1pt}}
% Subtema: título corto subrayado con filete.
\newcommand{\subtema}[1]{\par\vspace{3pt}{\bfseries\color{NavyBlue}#1}\par
  \nointerlineskip\vspace{1pt}{\color{NavyBlue!40}\rule{\linewidth}{0.4pt}}\par\vspace{1pt}}
% Una fórmula por renglón: fórmula a la izquierda, unidades a la derecha (gris).
\newcommand{\f}[2]{\par\(\displaystyle #1\)\hfill{\color{gray}\footnotesize #2}\par\vspace{2pt}}
\newcommand{\nota}[1]{\par{\color{gray}\footnotesize #1}\par\vspace{1pt}}

\begin{document}

{\large\bfseries Química General A1095 — Hoja de fórmulas}\hfill{\small UTN FRBA · Ing. Mecánica}\par
\vspace{1pt}{\color{NavyBlue}\rule{\linewidth}{0.8pt}}\par\vspace{2pt}

% multicols*: las columnas se llenan en orden, sin balancear (la hoja crece).
\begin{multicols*}{3}

% Título apenas más ancho que la columna: se comprime 6 % en horizontal para que entre en un renglón.
\capitulo{\scalebox{0.94}[1]{Cap. 4 — Reacciones en disolución acuosa}}

\subtema{Neutralización ácido–base}
\f{\text{ácido} + \text{base} \rightarrow \text{sal} + \text{H}_2\text{O}}{}
\f{\text{HCl} + \text{NaOH} \rightarrow \text{NaCl} + \text{H}_2\text{O}}{ej.}
\f{\text{H}^+ + \text{OH}^- \rightarrow \text{H}_2\text{O}}{ec. iónica neta}
{\footnotesize\setlength{\tabcolsep}{3pt}\begin{tabular}{@{}lll@{}}
 & \color{NavyBlue}\bfseries Ácido & \color{NavyBlue}\bfseries Base\\
Arrhenius & libera H\(^+\) (H\(_3\)O\(^+\)) & libera OH\(^-\)\\
Brønsted & dona H\(^+\) & acepta H\(^+\)\\
Lewis & acepta par e\(^-\) & dona par e\(^-\)\\
\end{tabular}}\par\vspace{2pt}

\subtema{RedOx {\mdseries\small(reducción–oxidación)}}
\f{\text{Zn} \rightarrow \text{Zn}^{2+} + 2\,\text{e}^-}{oxidación: \textbf{pierde} e\(^-\)}
\f{\text{Cu}^{2+} + 2\,\text{e}^- \rightarrow \text{Cu}}{reducción: \textbf{gana} e\(^-\)}
\nota{Reducir = bajar el n.º de ox.: suma e\(^-\) (negativos)}
\nota{Se oxida \(\rightarrow\) ag. reductor \quad se reduce \(\rightarrow\) ag. oxidante}

\capitulo{Cap. 5 — Gases}

\subtema{Gases ideales}
\f{P\,V = n\,R\,T}{\emph{Pibe Vago no Recibe Título}}
\nota{\(R = 0{,}0821\) L·atm/(mol·K) \(= 8{,}314\) J/(mol·K) \quad \(T\) [K] \(= t\) [°C] \(+ 273{,}15\)}

\capitulo{Cap. 6 — Termoquímica}

\subtema{Primera ley}
\f{\Delta E = q + w}{J}
\f{w = -P\,\Delta V}{L·atm \;(1\text{ L·atm} = 101{,}3\text{ J})}
\nota{1 cal = 4,184 J}
{\footnotesize\setlength{\tabcolsep}{2pt}\begin{tabular}{@{}ll@{}}
\(q>0\) absorbe (endo) & \(w>0\) sobre el sist. (\(\Delta V<0\))\\
\(q<0\) libera (exo) & \(w<0\) por el sist. (\(\Delta V>0\))
\end{tabular}}\par\vspace{2pt}

\subtema{Entalpía}
\f{H = E + PV}{J}
\f{\Delta H = \Delta E + P\,\Delta V}{J}
\f{q_p = \Delta H}{P = \text{cte}}
\f{\Delta E = \Delta H - RT\,\Delta n}{J \;(\(R = 8{,}314\) J/(mol·K))}
\nota{\(\Delta n\): mol de gas (prod − react) \quad \(T\) [K]}
\nota{\(\Delta H<0\) exotérmica \quad \(\Delta H>0\) endotérmica}

\subtema{Calorimetría}
\f{C = m\,s}{J/°C}
\nota{\(s\): calor específico [J/(g·°C)] \quad \(m\) [g]}
\f{q = m\,s\,\Delta t = C\,\Delta t}{J}
\nota{\(\Delta t = t_f - t_i\) [°C]}
{\footnotesize\setlength{\tabcolsep}{3pt}\begin{tabular}{@{}lr@{\hspace{10pt}}lr@{}}
\multicolumn{4}{@{}l}{\color{NavyBlue}\bfseries Calores específicos \(s\) [J/(g·°C)]}\\
Agua (H\(_2\)O líq.) & 4,184 & Hierro (Fe) & 0,444\\
Etanol & 2,46 & Cobre (Cu) & 0,385\\
Aluminio (Al) & 0,900 & Plomo (Pb) & 0,158\\
Grafito (C) & 0,720 & Mercurio (Hg) & 0,139\\
Diamante (C) & 0,502 & Oro (Au) & 0,129\\
\end{tabular}}\par\vspace{2pt}

\subtema{Entalpía de reacción}
\f{\Delta H^\circ_{\text{rxn}} = \textstyle\sum n\,\Delta H^\circ_f(\text{prod}) - \sum m\,\Delta H^\circ_f(\text{react})}{}
\nota{\(n, m\): coeficientes estequiométricos \hfill kJ/mol}
{\footnotesize\setlength{\tabcolsep}{3pt}\begin{tabular}{@{}lr@{\hspace{10pt}}lr@{}}
\multicolumn{4}{@{}l}{\color{NavyBlue}\bfseries \(\Delta H^\circ_f\) [kJ/mol] \mdseries(298 K, 1 atm)}\\
CO\(_{2}\)\,(g) & −393,5 & CH\(_{4}\)\,(g) & −74,8\\
H\(_{2}\)O\,(l) & −285,8 & C\(_{2}\)H\(_{2}\)\,(g) & +226,8\\
H\(_{2}\)O\,(g) & −241,8 & C\(_{3}\)H\(_{8}\)\,(g) & −103,8\\
CO\,(g) & −110,5 & C\(_{4}\)H\(_{10}\)\,(g) & −124,7\\
NH\(_{3}\)\,(g) & −46,2 & C\(_{8}\)H\(_{18}\)\,(l) & −269,8\\
NO\,(g) & +90,4 & C\(_{2}\)H\(_{5}\)OH\,(l) & −277,7\\
SO\(_{2}\)\,(g) & −296,9 & CaCO\(_{3}\)\,(s) & −1207,1\\
Fe\(_{2}\)O\(_{3}\)\,(s) & −822,2 & CaO\,(s) & −635,5\\
C diamante\,(s) & +1,9 &&\\
\end{tabular}}\par\vspace{1pt}
\nota{\(\Delta H^\circ_f = 0\): elemento puro estable (O\(_2\), H\(_2\), N\(_2\), C graf., Fe, Al)}
\f{\text{Hess: } \Delta H_{\text{total}} = \textstyle\sum \Delta H_{\text{etapas}}}{kJ}

\columnbreak
\capitulo{Cap. 11 — Fuerzas intermoleculares}

\nota{Dispersión < dip–dip < puente H < ion–dip}

\subtema{Puente de hidrógeno}
\f{\text{X–H}\cdots\text{:Y} \qquad \text{X, Y} = \text{N, O, F}}{}
\nota{H unido a N, O, F atraído por el par libre de otro N, O, F. Dipolo–dipolo especial, mucho más fuerte}
\nota{Da \(T_b\) y \(\Delta H_{\text{vap}}\) altos y menor \(P\) de vapor}
\nota{Hielo menos denso que el agua líq. (flota)}
{\footnotesize\setlength{\tabcolsep}{3pt}\begin{tabular}{@{}lrrr@{}}
\(T_b\) [°C] con p. H & H\(_2\)O\;100 & HF\;20 & NH\(_3\)\;−33\\
\hfill sin p. H & H\(_2\)S\;−60 & HCl\;−85 & PH\(_3\)\;−88\\
\end{tabular}}\par\vspace{2pt}

\subtema{Cambios de fase}
\f{\Delta H_{\text{sub}} = \Delta H_{\text{fus}} + \Delta H_{\text{vap}}}{kJ/mol}
\f{q = n\,\Delta H_{\text{fus}} \;(\text{o } \Delta H_{\text{vap}})}{kJ}
\nota{En el cambio de fase \(t\) = cte (meseta)}

\capitulo{Cap. 12 — Prop. físicas de las soluciones}

\subtema{Unidades de concentración}
\f{\%\,\text{masa} = \frac{m_{\text{soluto}}}{m_{\text{soluto}} + m_{\text{solvente}}}\times 100\%}{\%}
\f{X_A = \frac{n_A}{\sum n}}{adimensional}
\f{M = \frac{n_{\text{soluto}}}{V_{\text{solución}}}}{mol/L}
\f{m = \frac{n_{\text{soluto}}}{m_{\text{solvente}}}}{mol/kg}

\subtema{Propiedades coligativas}
\f{P_1 = X_1\,P_1^\circ}{atm, mmHg}
\nota{Raoult. 1 = solvente; \(P_1^\circ\): presión de vapor del solvente puro}
\f{\Delta T_b = K_b\,m}{°C}
\nota{\(\Delta T_b = T_b - T_b^\circ\) \quad \(K_b\) [°C/m]; agua: 0,52}
\f{\Delta T_f = K_f\,m}{°C}
\nota{\(\Delta T_f = T_f^\circ - T_f\) \quad \(K_f\) [°C/m]; agua: 1,86}
\f{\pi = M\,R\,T}{atm}
\nota{\(M\) [mol/L] \quad \(R = 0{,}0821\) L·atm/(mol·K) \quad \(T\) [K]}

\capitulo{Cap. 14 — Equilibrio químico}

\subtema{Ley de acción de masas}
\f{a\,A + b\,B \rightleftharpoons c\,C + d\,D}{}
\f{K_c = \frac{[C]^c\cdot[D]^d}{[A]^a\cdot[B]^b}}{[\;] en mol/L}
\f{K_P = \frac{P_C^{\,c}\cdot P_D^{\,d}}{P_A^{\,a}\cdot P_B^{\,b}}}{\(P\) en atm}
\f{K_P = K_c\,(R{\color{BrickRed}T})^{\Delta n}}{\(R = 0{,}0821\) L·atm/(mol·K)}
\nota{Solo gases: \(\Delta n = \sum\)coef prod \(-\sum\)coef react \quad \(T\) [K]}
\nota{\(K \gg 1\): predominan productos \quad \(K \ll 1\): reactivos}

\subtema{Cociente de reacción}
\f{Q_c = \frac{[C]_0^{\,c}\cdot[D]_0^{\,d}}{[A]_0^{\,a}\cdot[B]_0^{\,b}}}{[\;]\(_0\) iniciales}
{\footnotesize\setlength{\tabcolsep}{4pt}\begin{tabular}{@{}lll@{}}
\(Q_c < K_c\) & \(\rightarrow\) & avanza hacia productos\\
\(Q_c = K_c\) & \(\rightleftharpoons\) & en equilibrio\\
\(Q_c > K_c\) & \(\leftarrow\) & avanza hacia reactivos\\
\end{tabular}}\par\vspace{2pt}

\subtema{Le Châtelier}
\nota{Se desplaza para contrarrestar la perturbación}
{\footnotesize\setlength{\tabcolsep}{3pt}\begin{tabular}{@{}lll@{}}
\color{NavyBlue}\bfseries Perturbación & \color{NavyBlue}\bfseries Se desplaza a & \color{NavyBlue}\bfseries ¿\(K\)?\\
\(\uparrow\) react. o \(\downarrow\) prod. & productos & no\\
\(\uparrow P\) (\(\downarrow V\)) & menos mol de gas & no\\
\(\uparrow T\) endo (\(\Delta H>0\)) & productos & \textbf{sí}\\
\(\uparrow T\) exo (\(\Delta H<0\)) & reactivos & \textbf{sí}\\
Catalizador & no se desplaza & no\\
Gas inerte (\(V\) cte) & no se desplaza & no\\
\end{tabular}}\par\vspace{2pt}

\columnbreak
\capitulo{Cap. 15 — Ácidos y bases}

\subtema{Producto iónico del agua}
\f{2\,\text{H}_2\text{O} \rightleftharpoons \text{H}_3\text{O}^+ + \text{OH}^-}{}
\f{K_w = [\text{H}_3\text{O}^+]\cdot[\text{OH}^-] = [\text{H}^+]\cdot[\text{OH}^-]}{}
\f{K_w = 1{,}0\times10^{-14}}{a 25 °C}
{\footnotesize\setlength{\tabcolsep}{4pt}\begin{tabular}{@{}ll@{}}
Neutra & \([\text{H}^+] = [\text{OH}^-]\)\\
Ácida & \([\text{H}^+] > [\text{OH}^-]\)\\
Básica & \([\text{H}^+] < [\text{OH}^-]\)\\
\end{tabular}}\par\vspace{2pt}

\subtema{pH y pOH}
\f{\text{pH} = -\log[\text{H}^+] \;\Leftrightarrow\; [\text{H}^+] = 10^{-\text{pH}}}{}
\f{\text{pOH} = -\log[\text{OH}^-] \;\Leftrightarrow\; [\text{OH}^-] = 10^{-\text{pOH}}}{}
\f{\text{pH} + \text{pOH} = 14{,}00}{a 25 °C}

\subtema{Constantes de ionización}
\f{\text{HA} + \text{H}_2\text{O} \rightleftharpoons \text{H}_3\text{O}^+ + \text{A}^-}{ácido}
\f{K_a = \frac{[\text{H}^+]\cdot[\text{A}^-]}{[\text{HA}]}}{}
\f{\text{B} + \text{H}_2\text{O} \rightleftharpoons \text{BH}^+ + \text{OH}^-}{base}
\f{K_b = \frac{[\text{BH}^+]\cdot[\text{OH}^-]}{[\text{B}]}}{}
\nota{Mayor \(K\) \(\Rightarrow\) más fuerte \quad \([\;]\) de equilibrio}
\f{K_a\cdot K_b = K_w}{par conjugado (HA/A\(^-\), BH\(^+\)/B)}
\f{\%\ \text{ionización} = \frac{[\text{H}^+]_{\text{eq}}}{[\text{HA}]_0}\times 100\%}{\%}
\nota{\([\;]_0\): concentración inicial del ácido}

\capitulo{Cap. 18 — Energía libre y equilibrio}

\f{\Delta G^\circ = -R\,T\,\ln K}{J/mol}
\nota{\(R = 8{,}314\) J/(mol\(\cdot\)K) \quad \(T\) en K}
\nota{\(\Delta G^\circ < 0 \Leftrightarrow K > 1\) (favorece productos)}

\capitulo{Cap. 19 — Electroquímica}

\subtema{Celda galvánica}
{\footnotesize\setlength{\tabcolsep}{4pt}\begin{tabular}{@{}lll@{}}
\textbf{Á}nodo (−) & \textbf{O}xidación & vocales\\
\textbf{C}átodo (+) & \textbf{R}educción & consonantes\\
\end{tabular}}\par\vspace{1pt}
\nota{e\(^-\): ánodo \(\rightarrow\) cátodo por el cable}

\subtema{Potencial estándar de reducción}
\f{2\,\text{H}^+ + 2\,\text{e}^- \rightarrow \text{H}_2}{\(E^\circ = 0{,}00\) V (referencia)}
\nota{Mayor \(E^\circ_{\text{red}}\) \(\Rightarrow\) se reduce \(\Rightarrow\) \textbf{cátodo} (mejor oxidante)}
\f{E^\circ_{\text{celda}} = E^\circ_{\text{cátodo}} - E^\circ_{\text{ánodo}}}{fem [V]}
\nota{Estándar (°): 1 M, gases a 1 atm, 25 °C}
\nota{Ambos \(E^\circ_{\text{red}}\) de tabla, {\color{BrickRed}sin} invertir signo}
\f{\text{Ej.: } 0{,}34 {\color{BrickRed}\,-\,(-}0{,}76) = 1{,}10}{V}
\nota{Cátodo Cu\(^{2+}\)/Cu (\(+0{,}34\)) \quad ánodo Zn\(^{2+}\)/Zn (\(-0{,}76\))}

\subtema{Espontaneidad y trabajo eléctrico}
\f{\Delta G^\circ = -\underbrace{n\,F}_{q\ [\text{C}]}\,\underbrace{E^\circ_{\text{celda}}}_{[\text{V}]}}{C\(\cdot\)V = J}
\nota{\(n\): mol de e\(^-\) \quad \(F = 96\,500\) C/mol e\(^-\)}
\f{q = n\,F = I\,t}{C = A\(\cdot\)s}
\nota{\(I\): corriente [A = C/s] \quad \(t\): tiempo [s]}
\nota{fem = voltaje máx. \(\Rightarrow\) \(\Delta G\) = trabajo eléctrico máx.}
\nota{\(\Delta G\): energía libre, la disponible para hacer trabajo}
\nota{\(-\): el sistema hace trabajo sobre los alrededores}
\f{E^\circ_{\text{celda}} > 0 \;\Leftrightarrow\; \Delta G^\circ < 0 \;\Rightarrow\; \text{espontánea}}{}
\f{E^\circ_{\text{celda}} = \frac{R\,T}{n\,F}\,\ln K}{V}
\nota{\(R = 8{,}314\) J/(mol\(\cdot\)K) (no 0,0821) \quad \(T\) en K}

\subtema{Electrólisis {\mdseries\small(masa depositada)}}
\f{I\,t = q \overset{\div F}{\longrightarrow} \text{mol e}^- \overset{\div\,\text{e}^-\!/\text{ion}}{\longrightarrow} \text{mol} \overset{\times\mathcal{M}}{\longrightarrow} m}{g}
\nota{e\(^-\)/ion: de la hemirreacción (Cr\(^{3+}\) + 3 e\(^-\) \(\rightarrow\) Cr: 3)}

\end{multicols*}
\end{document}
